\section{Discussion}
\label{sec:discussion}

% src: Sections 5-8 (numbers as cited there); REPORT_SKELETON_VALIDATION.md (A1-A15); CLAIM_EVIDENCE_MAP.md C5 (common reading, test (a), consequences), C9
As shown in Sections~\ref{sec:results-gain}--\ref{sec:results-no-gain}, history lowers the error in every study, setting and model family but one. In this section, we discuss what history adds, when it helps and what this implies in practice, and then state the limitations of this work.

\noindent\textbf{What history adds.}
Section~\ref{subsec:history-dof} set out two ways in which a window can add to a single report: through the rates of change of the measurements, which require the reports in order, and through repeated noisy samples, which do not. In the macro cell, the gain comes from the second: the window averages out the error of the angle reports. On the controlled grid, the order of the reports also matters, so there the sequence models draw on the first mechanism as well.

\noindent\textbf{When history helps.}
The results fit one rule: history improves on whatever position information a single report already carries, and it cannot create information that is not there. With an angle, a single report locates the device in every study, and a window locates it better. Without an angle, a single report locates the device only where the model knows how the signal varies over the area for that device, as a fingerprinting system does once the area has been surveyed. The controlled grid and the $25\times25$ runs are of this kind: each model is trained on the first part of every walk and tested on the remainder, so it has in effect surveyed each device's radio map, and history sharpens its use of that map. Because the history gain is measured against the same model on the same split without history, the survey benefits both inputs equally. Without the survey, the same grid models fail~(Section~\ref{subsec:unseen-users}). The macro cell tests on users absent from training, so there a single-antenna report carries almost no usable information, and history adds nothing. The larger range of the macro cell does not explain this difference, since the RSS-history gain does not fall with range~(Section~\ref{subsec:hypothesis}). The one exception, single-antenna devices under mixed propagation, follows the same rule: there a window reveals whether the link is LOS or NLOS, something that no single report can show on its own.

\noindent\textbf{Practical case.}
For a deployment, the input matters more than the estimator. Once history is included, the trained models of the macro cell reach nearly the same error, and XGBoost, which trains in seconds, is among them. On the grid, a motion-model filter uses history better still, so the gain lies in the history itself and not in a particular learned model. Finally, the conditions that a deployment is most likely to vary, namely LOS propagation, a blocked line of sight and the interfering BSs, leave the gain intact.
